% Part: first-order-logic
% Chapter: tableaux
% Section: quantifier-rules

\documentclass[../../../include/open-logic-section]{subfiles}

\begin{document}

\olfileid{fol}{tab}{qrl}

\olsection{పరిమాణీకరణిక నియమాలు}

\subsection{$\lforall$కు నియమాలు}

\begin{defish}
\AxiomC{\sFmla{\True}{\lforall[x][!A(x)]}}
\RightLabel{\TRule{\True}{\forall}}
\UnaryInfC{\sFmla{\True}{!A(t)}}
\DisplayProof
\hfill
\AxiomC{\sFmla{\False}{\lforall[x][!A(x)]}}
\RightLabel{\TRule{\False}{\lforall}}
\UnaryInfC{\sFmla{\False}{!A(a)}}
\DisplayProof
\end{defish}

\TRule{\True}{\lforall}లో $t$ ఒక సంవృత పదం (అంటే చరరాశులు
లేని పదం). \TRule{\False}{\lforall}లో $a$ అనేది
\TRule{\False}{\lforall} నియమం పైనున్న శాఖలో ఎక్కడా కనిపించకూడని
ఒక !!a{constant}. \TRule{\False}{\forall} నిగమనం యొక్క
\emph{ఐగెన్ చరరాశి} అని $a$ను అంటాం.\footnote{పై నియమంలో $a$ ఒక
!!a{constant} అయినప్పటికీ “ఐగెన్ చరరాశి” అనే పేరునే ఉపయోగిస్తాం.
దీనికి చారిత్రక కారణాలున్నాయి.}

\subsection{$\lexists$కు నియమాలు}

\begin{defish}
\AxiomC{\sFmla{\True}{\lexists[x][!A(x)]}}
\RightLabel{\TRule{\True}{\lexists}}
\UnaryInfC{\sFmla{\True}{!A(a)}}
\DisplayProof
\hfill
\AxiomC{\sFmla{\False}{\lexists[x][!A(x)]}}
\RightLabel{\TRule{\False}{\lexists}}
\UnaryInfC{\sFmla{\False}{!A(t)}}
\DisplayProof
\end{defish}

ఇక్కడ కూడా $t$ ఒక సంవృత పదం; $a$ అనేది
\TRule{\True}{\lexists} నియమం పైనున్న శాఖలో కనిపించని ఒక
!!a{constant}. \TRule{\True}{\lexists} నిగమనం యొక్క
\emph{ఐగెన్ చరరాశి} అని $a$ను అంటాం.

\TRule{\False}{\lforall} లేదా \TRule{\True}{\lexists} నిగమనం
పైనున్న శాఖలో ఐగెన్ చరరాశి కనిపించకూడదనే షరతును
\emph{ఐగెన్ చరరాశి షరతు} అంటాం.

\begin{explain}
!!{variable}~$x$ స్వేచ్ఛగా ఉన్న ఒక !!a{formula}ను $!A$ సూచిస్తే,
దాన్ని $!A(x)$గా రాస్తామనే ఆచారాన్ని గుర్తుచేసుకోండి. అదే సందర్భంలో
$!A(t)$ అనేది~$\Subst{!A}{t}{x}$కు సంక్షిప్త రూపం. కాబట్టి
\TRule{\False}{\lexists} నియమాన్ని ఇలా కూడా రాయవచ్చు:
\begin{prooftree}
  \AxiomC{\sFmla{\False}{\lexists[x][!A]}}
  \RightLabel{\TRule{\False}{\lexists}}
  \UnaryInfC{\sFmla{\False}{\Subst{!A}{t}{x}}}
\end{prooftree}
$t$ ఇప్పటికే~$!A$లో కనిపించవచ్చని గమనించండి; ఉదాహరణకు $!A$
అనేది~$\Atom{\Obj P}{t,x}$ కావచ్చు. అందువల్ల
$\sFmla{\False}{\Atom{\Obj P}{t,t}}$ను
$ \sFmla{\False}{\lexists[x][\Atom{\Obj P}{t,x}]}$ నుంచి
నిగమించడం~\TRule{\False}{\lexists}ను సరిగ్గా వర్తింపజేయడమే. అయితే
\TRule{\False}{\lforall},~\TRule{\True}{\lexists}లలోని ఐగెన్
చరరాశి షరతులు !!{constant}~$a$ అనేది $!A$లో కనిపించకూడదని కోరుతాయి.
కాబట్టి
$\sFmla{\False}{\Atom{\Obj P}{a,a}}$ను
$\sFmla{\False}{\lforall[x][\Atom{\Obj P}{a,x}]}$ నుంచి
$\TRule{\False}{\lforall}$ ఉపయోగించి సరిగ్గా నిగమించలేం.
\end{explain}

\begin{explain}
\TRule{\True}{\lforall}, \TRule{\False}{\lexists}లలో పదం~$t$పై
ఎలాంటి పరిమితులూ లేవు. మరోవైపు, \TRule{\True}{\lexists},
\TRule{\False}{\lforall} నియమాల్లో ఐగెన్ చరరాశి షరతు,
!!{constant}~$a$ ఆయా నిగమనం పైనున్న శాఖల్లో ఎక్కడా కనిపించకూడదని
కోరుతుంది. వ్యవస్థ నిర్దుష్టంగా ఉండేలా ఈ షరతు అవసరం. ఈ షరతు లేకపోతే
$\lexists[x][\formula{A}(x)] \lif \lforall[x][\formula{A}(x)]$కు
కిందిది సంవృత !!{tableau} అవుతుంది:
\begin{center}
\begin{tableau}{}
  [\sFmla{\False}{\lexists[x][\formula{A}(x)] \lif \lforall[x][\formula{A}(x)]}, just=\TAss
    [\sFmla{\True}{\lexists[x][\formula{A}(x)]},
      just={\TRule{\False}{\lif}[1]}
      [\sFmla{\False}{\lforall[x][\formula{A}(x)]},
        just={\TRule{\False}{\lif}[1]}
        [\sFmla{\True}{\formula{A}(a)},
          just={\TRule{\True}{\lexists}[2]}
          [\sFmla{\False}{\formula{A}(a)},
            just={\TRule{\False}{\lforall}[3]}, close]
        ]
      ]
    ]
  ]
\end{tableau}
\end{center}
అయితే $\lexists[x][\formula{A}(x)] \lif
\lforall[x][\formula{A}(x)]$ చెల్లుబాటు కాదు.
\end{explain}

\end{document}
